\chapter*{TÓM TẮT KHÓA LUẬN}
\addcontentsline{toc}{chapter}{TÓM TẮT KHÓA LUẬN}

Trong lĩnh vực nhà hàng và dịch vụ ăn uống (F\&B), quy trình phục vụ tại bàn thường liên quan đến nhiều nhóm tác nhân như khách hàng, nhân viên phục vụ, bếp, thu ngân, chủ cửa hàng và bộ phận quản trị. Khi số lượng bàn, món ăn, đơn hàng và hình thức thanh toán tăng lên, cách vận hành thủ công hoặc các hệ thống bán hàng POS đơn lẻ dễ gặp khó khăn trong việc đồng bộ trạng thái đơn, kiểm soát tồn kho, cập nhật bếp, quản lý nhiều chi nhánh/đơn vị thuê bao và tích hợp thanh toán điện tử. Bên cạnh đó, mô hình đặt món qua mã QR giúp khách hàng thao tác trực tiếp trên thiết bị cá nhân, nhưng cũng đặt ra yêu cầu mới về xác thực phiên tại bàn, chống gửi trùng đơn, phối hợp với màn hình bếp KDS và bảo đảm dữ liệu của từng đơn vị thuê bao được cô lập.

Các sản phẩm thương mại trong lĩnh vực POS F\&B và đặt món qua mã QR cho thấy nhu cầu thực tế của bài toán, nhưng tài liệu công khai thường chỉ mô tả chức năng sản phẩm, không đủ để kết luận về ranh giới dịch vụ, quyền sở hữu dữ liệu, cơ chế nhất quán hay mức độ tách rời kiến trúc bên trong. Vì vậy, khóa luận không tập trung xây dựng hoặc đánh giá mô hình kinh doanh cho QRTable, cũng không nhằm phát minh lại POS hay đặt món qua mã QR. Thay vào đó, đề tài chọn bài toán SaaS POS tích hợp đặt món qua mã QR như một tình huống nghiên cứu kỹ thuật đủ phức tạp để trình bày một kiến trúc có ranh giới rõ, minh bạch và có thể kiểm chứng trong phạm vi báo cáo.

Từ định vị trên, khóa luận nghiên cứu và xây dựng QRTable, một nền tảng POS theo mô hình SaaS cho lĩnh vực F\&B, tích hợp đặt món qua mã QR và được thiết kế theo kiến trúc microservices. Hệ thống hướng đến các luồng vận hành chính gồm khởi tạo đơn vị thuê bao, quản lý thực đơn/bàn/mã QR, khách tham gia phiên và thao tác giỏ đặt món dùng chung, nhân viên xác nhận đơn, bếp tiếp nhận phiếu KDS, ghi nhận thanh toán tiền mặt hoặc VietQR/SePay, cùng các chức năng quản trị, phân quyền và báo cáo theo gói dịch vụ. Về kiến trúc, QRTable sử dụng BFF làm điểm vào cho các ứng dụng phía người dùng, tách các miền nghiệp vụ thành các dịch vụ như Catalog, Order, Kitchen, Payment, SaaS, User-Access và Authorizer, đồng thời dùng PostgreSQL/MongoDB theo ranh giới dịch vụ, Redis cho trạng thái khi vận hành, Kafka cho các sự kiện nghiệp vụ bất đồng bộ có chọn lọc, Keycloak/JWT cho nhân viên và quản trị, và cơ chế phiên QR riêng cho khách hàng.

Phương pháp thực hiện của khóa luận bắt đầu từ khảo sát bối cảnh POS, đặt món qua mã QR, SaaS, mô hình nhiều đơn vị thuê bao, kiến trúc microservices, kiến trúc hướng sự kiện, tính nhất quán dữ liệu và các yêu cầu bảo mật/phân quyền. Từ đó, đề tài phân tích tác nhân, ca sử dụng, yêu cầu chức năng và yêu cầu phi chức năng của QRTable; thiết kế ranh giới dịch vụ, quyền sở hữu dữ liệu, mô hình giao tiếp, chính sách cô lập đơn vị thuê bao, phân quyền theo vai trò và cơ chế xử lý lặp an toàn; sau đó triển khai các luồng nghiệp vụ đại diện trong một kho mã nguồn dùng chung theo Nx. Microservices giúp QRTable tách ranh giới nghiệp vụ và quyền sở hữu dữ liệu, nhưng đồng thời làm phát sinh các vấn đề của hệ thống phân tán như giao tiếp liên dịch vụ, nhất quán dữ liệu, phối hợp Saga, hợp đồng giữa dịch vụ và kiểm chứng kiến trúc. Phần đánh giá được tổ chức theo hướng đối chiếu yêu cầu, kiểm thử chức năng, kiểm chứng kiến trúc và phân tích ranh giới kết luận, thay vì chỉ dựa vào việc hệ thống chạy được trên một môi trường minh họa.

Kết quả đạt được là một hệ thống phần mềm có cấu trúc đầy đủ cho bài toán SaaS POS tích hợp đặt món qua mã QR. Các luồng quan trọng như phiên gọi món qua QR, giỏ đặt món dùng chung, gửi đơn, xác nhận đơn của nhân viên, trừ tồn kho, hàng đợi KDS, ghi nhận thanh toán, khởi tạo đơn vị thuê bao, bảng điều khiển và báo cáo theo quyền và phân quyền theo đơn vị thuê bao đã được mô tả, triển khai và đối chiếu bằng mã nguồn, kiểm thử, sơ đồ trình tự, bảng đối chiếu yêu cầu và các minh chứng tương ứng. Đặc biệt, khóa luận làm rõ cách QRTable áp dụng kiến trúc hướng sự kiện có chọn lọc: TCP/gRPC được dùng cho các yêu cầu cần phản hồi tức thời, Kafka dùng cho sự kiện nghiệp vụ bất đồng bộ, còn WebSocket chỉ đóng vai trò tín hiệu cập nhật/tải lại cho ứng dụng phía người dùng thay vì trở thành nguồn sự thật của trạng thái nghiệp vụ.

Trong phạm vi đánh giá, khóa luận không khẳng định hệ thống đã sẵn sàng vận hành công khai ở quy mô lớn, chưa đưa ra kết quả đo kiểm tải lớn và không xem các tích hợp với nhà cung cấp thực tế như SePay là cổng kiểm thử tự động mặc định. Thay vào đó, báo cáo xác định rõ mức độ minh chứng cho từng kết luận: phần nào đã được kiểm chứng bằng kiểm thử đơn vị, kiểm thử hợp đồng và kiểm thử tích hợp, phần nào được hỗ trợ bởi thiết kế và cấu trúc mã nguồn, và phần nào là hướng củng cố tiếp theo. Những hướng phát triển quan trọng gồm mở rộng kiểm thử tải, hoàn thiện minh chứng triển khai công khai, tăng cường kiểm chứng với nhà cung cấp thực tế, bổ sung cơ chế hàng đợi thao tác ngoại tuyến, nâng cấp khả năng quan sát hệ thống/củng cố vận hành và mở rộng các chức năng vận hành nâng cao cho nhà hàng.

\textbf{Từ khóa:} SaaS POS, đặt món qua mã QR, microservices, nhiều đơn vị thuê bao, kiến trúc hướng sự kiện, Kafka, Redis, Keycloak, VietQR, SePay.

\cleardoublepage
